\color{unimain}{\subsection{Токовая защита ошиновки}}
\color{black}

\begin{enumerate}

\item
Токовая защита ошиновки стороны автотрансформатора предназначена для защиты от коротких замыканий ошиновки отключенной стороны АТ, автоматически вводимая по факту отключенного положения коммутационных аппаратов защищаемой стороны. 

\item
Уставки ТЗО представлены в таблице \ref{app_set:tzo}.

\item
В составе ТЗО предусмотрены:
\begin{itemize}
\item ступень, реагирующая на величину фазных токов (ТЗО-Ф);
\item ступень, реагирующая на величину тока нулевой последовательности (ТЗО-НП).
\end{itemize}

\item
На рисунке \ref{pict_tzo:TZO} приведен алгоритм работы ТЗО. 

\begin{figure}[H]
\centering
\includegraphics[width=1\textwidth,height=1\textheight,keepaspectratio]{\fbpath/07.02 РЗАТ/041.0191 ТЗО/img/TZO.pdf}
\captionof{figure}{Функциональная схема логики ТЗО}\label{pict_tzo:TZO}
\end{figure}

\item
Ввод в работу ступеней осуществляется с помощью программных переключателей \textbf{«ТЗО-Ф»} и \textbf{«ТЗО-НП»}. Пороги срабатывания токов ТЗО-Ф и ТЗО-НП задаются с помощью уставок \textbf{«Iср ТЗО-Ф»} и \textbf{«3I0ср ТЗО-НП»} соответственно, а выдержки времени срабатывания ступеней регулируются уставками \textbf{«Tср ТЗО-Ф»} и \textbf{«Tср ТЗО-НП»}. 

\item
Оперативный ввод/вывод ТЗО производится с помощью БУ (блоков управления функциями) \textit{«Вывод ТЗО»}. 

\item
В зависимости от схемы подключения АТ имеется возможность выбрать условия пуска автоматического ввода ТЗО, задается программным переключателем \textbf{«Авт ввод ТЗО»}, имеющими три положения:
\begin{itemize}
\item \textit{«Р»} -- Пуск ТЗО по факту отключенного положения ТР (Р1) и ОТР (Р2);
\item \textit{«В»} -- Пуск ТЗО по факту отключенного положения В (В1) и ОВ (В2);
\item \textit{«Р, В»} -- Пуск ТЗО по факту отключенного положения ТР (Р1) и ОТР (Р2) или В (В1) и ОВ (В2).
\end{itemize}

\item
Сигналы отключенного положения разъединителей назначаются на конфигурируемые входа \textit{«ТР отключен»} и \textit{«ОТР отключен»}.

\item
Предусмотрена также возможность оперативного ввода ТЗО (независимо от положения КА) с помощью команды управления \textit{«Опер. ввод ТЗО»}.

\item
Ввод контроля положения обходного трансформаторного разъединителя задается программным переключателем \textbf{«Контроль ОТР»}, а ввод контроля положения ОВ (В2) задается программным переключателем \textbf{«Контроль В2»}.

\item
Срабатывание ТЗО предусматривает действие на отключение АТ со всех смежных сторон с запретом АПВ.

\end{enumerate} 